\documentclass{article}
\usepackage[utf8]{inputenc}
\usepackage[T1]{fontenc}
\usepackage{graphicx}
\usepackage{hyperref}
\usepackage{geometry}
\geometry{margin=1in}

\title{ARCHITECTURAL ORDERS and COLUMNS (Plate XXIX)}
\date{Encyclopaedia Britannica (1771)}

\begin{document}
\maketitle

\section*{Historical source}
\textbf{Entry ID:} \texttt{entry-britannica-volume-191253799-row-374-374} \\
\textbf{Volume ID:} 191253799 \\
\textbf{Character Count:} 966

\subsection*{Extracted Text}
[HISTORICAL COPPERPLATE ENGRAVING: ARCHITECTURAL ORDERS and COLUMNS (Plate XXIX)]
Category: Architecture and Classical Design
Visual Description: Copperplate engraving displaying the Five Classical Orders of Architecture: Tuscan, Doric, Ionic, Corinthian, and Composite. Details pedestal, base, shaft, capital, architrave, frieze, and cornice proportions.

[Nemotron Super VL Analysis]: Based on the visual evidence, the provided image displays a blank, aged page with minor staining (foxing) and a binding edge visible on the left. There are no visible headings, figures, columns, or engravings present in the image to transcribe or describe.
Identified Engraved Figures:
  * Fig. 1: Tuscan and Doric orders with entablature proportions
  * Fig. 2: Ionic order showing spiral volute geometry
  * Fig. 3: Corinthian and Composite capitals decorated with acanthus leaves
Related Scientific Subjects: Classical order, Doric order, Ionic order, Corinthian order, Column

\subsection*{Matched Wikipedia Article}
\textbf{Title:} History of the Encyclopædia Britannica \\
\textbf{Page ID:} 9004976 \\
\textbf{Revision ID:} 1340071003 \\
\textbf{URL:} \url{https://en.wikipedia.org/wiki/History\_of\_the\_Encyclop\%C3\%A6dia\_Britannica}


\section*{Source record}
\begin{itemize}
\item \textbf{Dataset:} \texttt{NationalLibraryOfScotland/encyclopaedia\_britannica\_illustrated}
\item \textbf{Dataset Revision:} \texttt{309f765b1ecc83e0a7f7762e44b4a927a7555443}
\item \textbf{Page ID:} \texttt{britannica:309f765b1ecc83e0a7f7762e44b4a927a7555443:train:374} \\
  \textbf{Filename:} \texttt{page-0374-2b13c7e2a0e0.jpg} \\
  \textbf{Row Index:} 374 \\
  \textbf{SHA256:} \texttt{2b13c7e2a0e0431e...}
\end{itemize}

\begin{figure}[h]
\centering
\includegraphics[width=0.85\textwidth]{../scans/page-0374-2b13c7e2a0e0.jpg}
\caption*{Source scan: \texttt{page-0374-2b13c7e2a0e0.jpg}}
\end{figure}

\end{document}
